% GENERATED FILE. Edit canonical lessons, then run npm run generate:books.
% canonical-source-hash: fnv1a64:73d4fb4a6fc1bd8a
% canonical-lessons: UR-C152-bimari, UR-C152-zukam, UR-C152-khansi, UR-C152-sar-dard, UR-C152-pet-dard

\chapter{Illness, Head cold, Cough, Headache, Stomach ache}
\label{ch:ur-illness-head-cold-cough-headache-stomach-ache}
\hlchaptermodality{152}

\begin{chapteropening}
\textbf{By the end of this chapter:} \emph{I can say an illness, a head cold, a cough, a headache and a stomach ache.}

Complete the last lesson of chapter 152: I can say an illness, a head cold, a cough, a headache and a stomach ache.
\end{chapteropening}

\section[\texorpdfstring{\ur{بیماری} (bīmārī)}{بیماری (bīmārī)}]{\ur{بیماری} (bīmārī) — an illness}

\label{lesson:UR-C152-bimari}



\begin{quote}
Before the new one: say the Urdu for to cough, then the Urdu for to breathe.
\end{quote}

\subsection*{You'll want to know: \ur{بیماری}}
\textbf{\ur{بیماری}} — \emph{bīmārī} — \textquotedblleft{}an illness\textquotedblright{}.

\begin{cousinweb}
\textbf{\ur{بیماری}} (\emph{bīmārī}) is \emph{bīmār}, \textquotedblleft{}ill, sick\textquotedblright{}, usually traced to Persian \emph{bīmār}, with the ending \emph{-ī} that makes a noun.
\end{cousinweb}

\begin{grammarlens}[title={What you've built}]
One more word for this chapter.
\end{grammarlens}

\subsection*{Your turn}
\begin{itemize}
\raggedright
  \item \emph{Say it:} \emph{bīmārī}
  \item \emph{Say it:} \emph{bīmārī}, once more
  \item \emph{Say it:} say \emph{bīmārī}
  \item \emph{Recall:} say the Urdu for to cough, then the Urdu for to breathe, then say \emph{bīmārī} again
\end{itemize}

\begin{tcolorbox}[breakable,colback=teal!4,colframe=teal!35!black,title={Before you move on}]
What does \ur{بیماری} mean? (\textquotedblleft{}An illness\textquotedblright{}.) Say it once more.
\end{tcolorbox}
\section[\texorpdfstring{\ur{زکام} (zukām)}{زکام (zukām)}]{\ur{زکام} (zukām) — a head cold}

\label{lesson:UR-C152-zukam}



\begin{quote}
Before the new one: say the Urdu for to breathe, then the Urdu for an illness.
\end{quote}

\subsection*{You'll want to know: \ur{زکام}}
\textbf{\ur{زکام}} — \emph{zukām} — \textquotedblleft{}a head cold\textquotedblright{}.

\begin{cousinweb}
\textbf{\ur{زکام}} (\emph{zukām}) is a loan from Arabic \emph{zukām}, \textquotedblleft{}a cold, catarrh\textquotedblright{}, which reached Urdu through Persian. Its Arabic root is \emph{z-k-m} (Lane).
\end{cousinweb}

\begin{grammarlens}[title={What you've built}]
One more word for this chapter.
\end{grammarlens}

\subsection*{Your turn}
\begin{itemize}
\raggedright
  \item \emph{Say it:} \emph{zukām}
  \item \emph{Say it:} \emph{zukām}, once more
  \item \emph{Say it:} say \emph{zukām}
  \item \emph{Recall:} say the Urdu for to breathe, then the Urdu for an illness, then say \emph{zukām} again
\end{itemize}

\begin{tcolorbox}[breakable,colback=teal!4,colframe=teal!35!black,title={Before you move on}]
What does \ur{زکام} mean? (\textquotedblleft{}A head cold\textquotedblright{}.) Say it once more.
\end{tcolorbox}
\section[\texorpdfstring{\ur{کھانسی} (khāṅsī)}{کھانسی (khāṅsī)}]{\ur{کھانسی} (khāṅsī) — a cough}

\label{lesson:UR-C152-khansi}



\begin{quote}
Before the new one: say the Urdu for an illness, then the Urdu for a head cold.
\end{quote}

\subsection*{You'll want to know: \ur{کھانسی}}
\textbf{\ur{کھانسی}} — \emph{khāṅsī} — \textquotedblleft{}a cough\textquotedblright{}.

\begin{cousinweb}
\textbf{\ur{کھانسی}} (\emph{khāṅsī}) is a tadbhava, a word inherited from Sanskrit and reshaped by sound change. It comes from \emph{kāsa}, \textquotedblleft{}cough\textquotedblright{} (CDIAL 3138). CDIAL notes the \emph{kh-} as imitating the sound.
\end{cousinweb}

\begin{grammarlens}[title={What you've built}]
One more word for this chapter.
\end{grammarlens}

\subsection*{Your turn}
\begin{itemize}
\raggedright
  \item \emph{Say it:} \emph{khāṅsī}
  \item \emph{Say it:} \emph{khāṅsī}, once more
  \item \emph{Say it:} say \emph{khāṅsī}
  \item \emph{Recall:} say the Urdu for an illness, then the Urdu for a head cold, then say \emph{khāṅsī} again
\end{itemize}

\begin{tcolorbox}[breakable,colback=teal!4,colframe=teal!35!black,title={Before you move on}]
What does \ur{کھانسی} mean? (\textquotedblleft{}A cough\textquotedblright{}.) Say it once more.
\end{tcolorbox}
\section[\texorpdfstring{\ur{سر} \ur{درد} (sar dard)}{سر درد (sar dard)}]{\ur{سر} \ur{درد} (sar dard) — a headache}

\label{lesson:UR-C152-sar-dard}



\begin{quote}
Before the new one: say the Urdu for a head cold, then the Urdu for a cough.
\end{quote}

\subsection*{You'll want to know: \ur{سر} \ur{درد}}
\textbf{\ur{سر} \ur{درد}} — \emph{sar dard} — \textquotedblleft{}a headache\textquotedblright{}.

\begin{cousinweb}
\textbf{\ur{سر} \ur{درد}} (\emph{sar dard}) joins \emph{sar}, \textquotedblleft{}head\textquotedblright{}, and \emph{dard}, \textquotedblleft{}pain\textquotedblright{}, both usually traced to Persian.
\end{cousinweb}

\begin{grammarlens}[title={What you've built}]
One more word for this chapter.
\end{grammarlens}

\subsection*{Your turn}
\begin{itemize}
\raggedright
  \item \emph{Say it:} \emph{sar dard}
  \item \emph{Say it:} \emph{sar dard}, once more
  \item \emph{Say it:} say \emph{sar dard}
  \item \emph{Recall:} say the Urdu for a head cold, then the Urdu for a cough, then say \emph{sar dard} again
\end{itemize}

\begin{tcolorbox}[breakable,colback=teal!4,colframe=teal!35!black,title={Before you move on}]
What does \ur{سر} \ur{درد} mean? (\textquotedblleft{}A headache\textquotedblright{}.) Say it once more.
\end{tcolorbox}
\section[\texorpdfstring{\ur{پیٹ} \ur{درد} (peṭ dard)}{پیٹ درد (peṭ dard)}]{\ur{پیٹ} \ur{درد} (peṭ dard) — a stomach ache}

\label{lesson:UR-C152-pet-dard}



\begin{quote}
Before the new one: say the Urdu for a cough, then the Urdu for a headache.
\end{quote}

\subsection*{You'll want to know: \ur{پیٹ} \ur{درد}}
\textbf{\ur{پیٹ} \ur{درد}} — \emph{peṭ dard} — \textquotedblleft{}a stomach ache\textquotedblright{}.

\begin{cousinweb}
\textbf{\ur{پیٹ} \ur{درد}} (\emph{peṭ dard}) joins \emph{peṭ}, \textquotedblleft{}belly\textquotedblright{}, and \emph{dard}, \textquotedblleft{}pain\textquotedblright{}. CDIAL lists \emph{peṭ} under a reconstructed \emph{*peṭṭa}, \textquotedblleft{}belly\textquotedblright{} (CDIAL 8376), and notes that it may be the same word as forms meaning \textquotedblleft{}basket\textquotedblright{}; \emph{dard} is usually traced to Persian.
\end{cousinweb}

\begin{grammarlens}[title={What you've built}]
One more word for this chapter.
\end{grammarlens}

\subsection*{Your turn}
\begin{itemize}
\raggedright
  \item \emph{Say it:} \emph{peṭ dard}
  \item \emph{Say it:} \emph{peṭ dard}, once more
  \item \emph{Say it:} say \emph{peṭ dard}
  \item \emph{Recall:} say the Urdu for a cough, then the Urdu for a headache, then say \emph{peṭ dard} again
\end{itemize}

\begin{tcolorbox}[breakable,colback=teal!4,colframe=teal!35!black,title={Before you move on}]
What does \ur{پیٹ} \ur{درد} mean? (\textquotedblleft{}A stomach ache\textquotedblright{}.) Say it once more.
\end{tcolorbox}
